% GENERATED FILE. Edit canonical lessons, then run npm run generate:books.
% canonical-source-hash: fnv1a64:928484ee7330e6a3
% canonical-lessons: ES-C750-vinculo, ES-C750-mimar, ES-C750-convivir, ES-C750-colega, ES-C750-apoyo

\chapter{Bajo el mismo techo}
\label{ch:es-bajo-el-mismo-techo}
\hlchaptermodality{750}

\begin{chapteropening}
\textbf{By the end of this chapter:} \emph{I can talk about sharing a home, pampering each other and the bonds that support us.}

Complete the last lesson of chapter 750: I can talk about sharing a home, pampering each other and the bonds that support us.
\end{chapteropening}

\section[el vínculo]{el vínculo — the bond, the link}

\label{lesson:ES-C750-vinculo}



\begin{quote}
Before the new one: say the Spanish for the great-grandfather, then the Spanish for the orphan.
\end{quote}

\subsection*{You'll want to know: el vínculo}
\textbf{el vínculo} — \textquotedblleft{}the bond, the link\textquotedblright{}.

\begin{quote}
Tengo un vínculo especial con mi abuela.
\end{quote}

\begin{quote}
\emph{I have a special bond with my grandmother.}
\end{quote}

\begin{cousinweb}
\emph{vínculo} is a learned borrowing from Latin \emph{vinculum} \textquotedblleft{}bond, link\textquotedblright{} (EtymDB), built on \emph{vincīre} \textquotedblleft{}to bind\textquotedblright{}.
\end{cousinweb}

\begin{grammarlens}[title={What you've built}]
One more word for telling, explaining and describing at B1.
\end{grammarlens}

\subsection*{Your turn}
\begin{itemize}
\raggedright
  \item \emph{Say it:} \emph{el vínculo}
  \item \emph{Say it:} \emph{el vínculo}, once more
  \item \emph{Say it:} a sentence of your own with \emph{el vínculo}
  \item \emph{Recall:} say the Spanish for the great-grandfather, then the Spanish for the orphan, then say \emph{el vínculo} again
\end{itemize}

\begin{tcolorbox}[breakable,colback=teal!4,colframe=teal!35!black,title={Before you move on}]
What does el vínculo mean? (\textquotedblleft{}The bond, the link\textquotedblright{}.) Say it once more.
\end{tcolorbox}
\section[mimar]{mimar — to spoil, to pamper}

\label{lesson:ES-C750-mimar}



\begin{quote}
Before the new one: say the Spanish for the orphan, then the Spanish for the bond.
\end{quote}

\subsection*{You'll want to know: mimar}
\textbf{mimar} — \textquotedblleft{}to spoil, to pamper\textquotedblright{}.

\begin{quote}
Los abuelos miman mucho a los nietos.
\end{quote}

\begin{quote}
\emph{Grandparents really spoil their grandchildren.}
\end{quote}

\emph{Mimado} — spoilt. \emph{Los mimos} — cuddles, pampering.

\begin{cousinweb}
\emph{mimar} is built on \emph{mimo} \textquotedblleft{}a cuddle, pampering\textquotedblright{}, which is usually traced to Latin \emph{mīmus} \textquotedblleft{}mime, clown\textquotedblright{}, from Greek \emph{mîmos}: the playful faces made to a child.
\end{cousinweb}

\begin{grammarlens}[title={What you've built}]
One more word for telling, explaining and describing at B1.
\end{grammarlens}

\subsection*{Your turn}
\begin{itemize}
\raggedright
  \item \emph{Say it:} \emph{mimar}
  \item \emph{Say it:} \emph{mimar}, once more
  \item \emph{Say it:} a sentence of your own with \emph{mimar}
  \item \emph{Recall:} say the Spanish for the orphan, then the Spanish for the bond, then say \emph{mimar} again
\end{itemize}

\begin{tcolorbox}[breakable,colback=teal!4,colframe=teal!35!black,title={Before you move on}]
What does mimar mean? (\textquotedblleft{}To spoil, to pamper\textquotedblright{}.) Say it once more.
\end{tcolorbox}
\section[convivir]{convivir — to live together}

\label{lesson:ES-C750-convivir}



\begin{quote}
Before the new one: say the Spanish for the bond, then the Spanish for to spoil.
\end{quote}

\subsection*{You'll want to know: convivir}
\textbf{convivir} — \textquotedblleft{}to live together\textquotedblright{}.

\begin{quote}
Convivo con dos compañeros de piso.
\end{quote}

\begin{quote}
\emph{I live with two flatmates.}
\end{quote}

\emph{La convivencia} — living together.

\begin{cousinweb}
\emph{convivir} comes from Latin \emph{convīvere} \textquotedblleft{}to live together\textquotedblright{} (EtymDB), built on \emph{vīvere} \textquotedblleft{}to live\textquotedblright{}.
\end{cousinweb}

\begin{grammarlens}[title={What you've built}]
One more word for telling, explaining and describing at B1.
\end{grammarlens}

\subsection*{Your turn}
\begin{itemize}
\raggedright
  \item \emph{Say it:} \emph{convivir}
  \item \emph{Say it:} \emph{convivir}, once more
  \item \emph{Say it:} a sentence of your own with \emph{convivir}
  \item \emph{Recall:} say the Spanish for the bond, then the Spanish for to spoil, then say \emph{convivir} again
\end{itemize}

\begin{tcolorbox}[breakable,colback=teal!4,colframe=teal!35!black,title={Before you move on}]
What does convivir mean? (\textquotedblleft{}To live together\textquotedblright{}.) Say it once more.
\end{tcolorbox}
\section[el colega]{el colega — the colleague; mate}

\label{lesson:ES-C750-colega}



\begin{quote}
Before the new one: say the Spanish for to spoil, then the Spanish for to live together.
\end{quote}

\subsection*{You'll want to know: el colega}
\textbf{el colega} — \textquotedblleft{}the colleague; mate\textquotedblright{}.

\begin{quote}
Voy a cenar con unos colegas.
\end{quote}

\begin{quote}
\emph{I am having dinner with some mates.}
\end{quote}

At work it is a colleague; in Spain, colloquially, a mate.

\begin{cousinweb}
\emph{colega} is a learned borrowing from Latin \emph{collēga} \textquotedblleft{}a partner in office\textquotedblright{} (EtymDB).
\end{cousinweb}

\begin{grammarlens}[title={What you've built}]
One more word for telling, explaining and describing at B1.
\end{grammarlens}

\subsection*{Your turn}
\begin{itemize}
\raggedright
  \item \emph{Say it:} \emph{el colega}
  \item \emph{Say it:} \emph{el colega}, once more
  \item \emph{Say it:} a sentence of your own with \emph{el colega}
  \item \emph{Recall:} say the Spanish for to spoil, then the Spanish for to live together, then say \emph{el colega} again
\end{itemize}

\begin{tcolorbox}[breakable,colback=teal!4,colframe=teal!35!black,title={Before you move on}]
What does el colega mean? (\textquotedblleft{}The colleague; mate\textquotedblright{}.) Say it once more.
\end{tcolorbox}
\section[el apoyo]{el apoyo — support}

\label{lesson:ES-C750-apoyo}



\begin{quote}
Before the new one: say the Spanish for to live together, then the Spanish for the colleague.
\end{quote}

\subsection*{You'll want to know: el apoyo}
\textbf{el apoyo} — \textquotedblleft{}support\textquotedblright{}.

\begin{quote}
Gracias por tu apoyo.
\end{quote}

\begin{quote}
\emph{Thank you for your support.}
\end{quote}

\emph{Apoyar} — to support.

\begin{cousinweb}
\emph{apoyo} is built on \emph{apoyar} (EtymDB), which is usually traced to Italian \emph{appoggiare} \textquotedblleft{}to lean on\textquotedblright{}, from Latin \emph{podium} \textquotedblleft{}a raised base\textquotedblright{}.
\end{cousinweb}

\begin{grammarlens}[title={What you've built}]
One more word for telling, explaining and describing at B1.
\end{grammarlens}

\subsection*{Your turn}
\begin{itemize}
\raggedright
  \item \emph{Say it:} \emph{el apoyo}
  \item \emph{Say it:} \emph{el apoyo}, once more
  \item \emph{Say it:} a sentence of your own with \emph{el apoyo}
  \item \emph{Recall:} say the Spanish for to live together, then the Spanish for the colleague, then say \emph{el apoyo} again
\end{itemize}

\begin{tcolorbox}[breakable,colback=teal!4,colframe=teal!35!black,title={Before you move on}]
What does el apoyo mean? (\textquotedblleft{}Support\textquotedblright{}.) Say it once more.
\end{tcolorbox}
